Wir betrachten auf dem $\R^2$ das stetige Vektorfeld $F$, das durch

\mavergleichskettedisp
{\vergleichskette
{F(x,y)
}
{ =} { { \frac{   1 }{  x^2+y^2 } }  e_1
}
{ } {
}
{ } {
}
{ } {
}
}
{}{}{}
gegeben sei. Dieses hat keine Nullstelle und ist stetig. Wir transportieren dieses Vektorfeld mittels der stereographischen Projektion auf

\mavergleichskette
{\vergleichskette
{  \R^2
}
{ \cong }{ S^2 \setminus \{N\}
}
{  }{
}
{    }{
}
{   }{
}
}
{}{}{}
und ergänzen es im Nordpol durch den Wert

\mavergleichskette
{\vergleichskette
{ 0
}
{ \in }{  T_NS^2
}
{  }{
}
{    }{
}
{   }{
}
}
{}{}{.}
Wir behaupten, dass dieses Vektorfeld stetig ist. Dazu sei
\mathl{P_n}{} sei eine Folge auf $S^2$, die gegen $N$ konvergiert. Dabei können wir direkt annehmen, dass alle

\mavergleichskette
{\vergleichskette
{P_n
}
{ \neq }{N
}
{  }{
}
{    }{
}
{   }{
}
}
{}{}{}
sind. Das Kartenbild dieser Folge ist

\mavergleichskettedisp
{\vergleichskette
{P'_n
}
{ =} { (x_n,y_n)
}
{ } {
}
{ } {
}
{ } {
}
}
{}{}{,}
und da $P_n$ gegen den Nordpol konvergiert, divergiert
\mathl{\betrag { P_n' }}{} bestimmt gegen $\infty$. Daher konvergiert die Folge

\mavergleichskettedisp
{\vergleichskette
{ F(P'_n)
}
{ =} { F(x_n , y_n)
}
{ =} { { \frac{   1 }{  x^2+y^2 } } e_1
}
{ } {
}
{ } {
}
}
{}{}{}
gegen $0$.

 [[Kategorie:Latexseite]]
